\documentclass[12pt, a4paper]{article}

\usepackage{geometry}
\geometry{margin=1.15in, top=1.2in, bottom=1.2in}

\usepackage{fontspec}
\setmainfont{TeX Gyre Pagella}
\usepackage{microtype}

\usepackage{amsmath, amsthm, amssymb, mathtools}
\usepackage{unicode-math}
\setmathfont{TeX Gyre Pagella Math}

\usepackage{hyperref}
\hypersetup{colorlinks=true, linkcolor=black, citecolor=black, urlcolor=black}

\usepackage{booktabs, enumitem, xcolor, setspace}
\usepackage{abstract}

% ── Theorem environments ─────────────────────────────────────────────────────
\theoremstyle{plain}
\newtheorem{theorem}{Theorem}[section]
\newtheorem{proposition}[theorem]{Proposition}
\newtheorem{lemma}[theorem]{Lemma}
\newtheorem{corollary}[theorem]{Corollary}

\theoremstyle{definition}
\newtheorem{definition}[theorem]{Definition}
\newtheorem{example}[theorem]{Example}

\theoremstyle{remark}
\newtheorem{remark}[theorem]{Remark}

% ── Notation ─────────────────────────────────────────────────────────────────
\newcommand{\calX}{\mathcal{X}}
\newcommand{\calM}{\mathcal{M}}
\newcommand{\calF}{\mathcal{F}}
\newcommand{\calC}{\mathcal{C}}
\newcommand{\calB}{\mathcal{B}}
\newcommand{\calO}{\mathcal{O}}
\newcommand{\calR}{\mathcal{R}}
\newcommand{\bR}{\mathbb{R}}
\newcommand{\vpi}{\pi}
\newcommand{\ve}{\varepsilon}
\newcommand{\dM}{d_{\calM}}
\newcommand{\norm}[1]{\left\lVert #1 \right\rVert}
\newcommand{\inner}[2]{\langle #1,\, #2 \rangle}
\newcommand{\fiber}[1]{\vpi^{-1}(#1)}

\renewcommand{\abstractnamefont}{\normalfont\small\bfseries}
\renewcommand{\abstracttextfont}{\normalfont\small}

% ── Title ────────────────────────────────────────────────────────────────────
\title{\textbf{Document Perturbation as Constrained Navigation\\
on a Semantic Manifold}}
\author{Flyxion\\
\small Independent Researcher}
\date{}

% ════════════════════════════════════════════════════════════════════════════
\begin{document}
\maketitle
\thispagestyle{empty}

% ── Abstract ─────────────────────────────────────────────────────────────────
\begin{abstract}
\noindent
We introduce a framework for generating controlled document variations through
constrained navigation in semantic space.  Traditional document augmentation
relies on synonym replacement, paraphrasing, or stochastic language-model
sampling, providing little formal control over semantic drift.  We propose
a \emph{semantic-budget framework} in which document transformations are
perturbation operators acting on a semantic manifold $\calM$.  A document is
a point in an embedding space; admissible transformations are constrained by
a finite semantic distance budget, local curvature bounds, and conservation
laws preserving core concepts.  The framework is developed as exploration of
the fibers of a projection $\vpi : \calX \to \calM$, where $\calX$ is the
space of syntactically valid textual realizations and $\calM$ encodes semantic
identity.  Motivated by colored-noise diffusion sampling
\cite{davidson2026colorednoisediffusionsampling}, we introduce a
\emph{semantic completion field} $\gamma(c, t) \in [0,1]$ and prove that
budget allocation weighted by $1 - \gamma(c, t)$ maximizes semantic diversity
subject to a finite perturbation budget, constituting semantic colored noise.
Mutation profiles are modelled as profile-dependent metric tensors on $\calM$,
enabling academic, legal, journalistic, technical, and literary registers to
define semantic equivalence according to their own geometries.  Applications
include data augmentation, retrieval robustness evaluation, deduplication
benchmarking, and adversarial robustness testing.
\end{abstract}

\newpage
\tableofcontents
\newpage

% ════════════════════════════════════════════════════════════════════════════
\section{Introduction}
\label{sec:intro}
% ════════════════════════════════════════════════════════════════════════════

Controlled variation of natural-language documents underlies data augmentation
for low-resource learning, robustness evaluation for retrieval and
question-answering systems, deduplication benchmark construction, plagiarism
detection calibration, and synthetic parallel corpus generation.  Despite the
centrality of these applications, the dominant approach remains conceptually
simple: select words at random, substitute synonyms found in a lexical database,
and hope that the resulting text is both meaningful and usefully different from
the original.

The limitations of this approach are well understood.  WordNet-based substitution
\cite{miller1995wordnet} treats all synonyms as equiprobable and ignores
contextual fit, producing phrases such as \emph{the investigator cautiously
promulgated the constitution of the antique edifice} as a paraphrase of
\emph{the scientist carefully published the structure of the ancient building}.
Large-language-model paraphrasing \cite{brown2020gpt3} resolves the contextual
problem but introduces a different one: with no formal constraint on semantic
drift, an unconstrained model may produce a document that preserves surface
fluency while substantially altering meaning, logical structure, or
domain-specific claims.  Neither approach provides a principled account of
\emph{how much} the resulting document has changed or \emph{which aspects} of
the original have been preserved.

Most NLP papers treat the problem as follows: a document is a string; we
transform the string; we measure whether the transformation was successful.
The ontological status of the document is never questioned.  We take a
different position.  A document is a point on a semantic manifold; the text
file is merely one coordinate representation of an underlying semantic object.
Two different strings---\emph{The scientist examined the structure} and
\emph{The researcher analyzed the architecture}---may occupy nearly the same
region of the manifold.  The geometric object is primary; the text is secondary.

This reframing elevates document perturbation from an NLP engineering task to a
problem in constrained trajectory planning on a manifold.  The framework admits
a natural formal language: fiber bundles, metric tensors, holonomic constraints,
geodesic curvature, and completion-weighted energy allocation.  Our central
contribution is to show that this language does genuine mathematical work,
enabling results that are unavailable within the string-substitution paradigm.

We make the following contributions.

\begin{enumerate}[label=(\roman*)]
\item We introduce \emph{semantic equivalence classes} via an approximate
  equivalence relation on the manifold and show that fibers are exactly these
  classes (\S\ref{sec:manifold}).
\item We develop the \emph{semantic fiber framework} (\S\ref{sec:fiber}),
  defining reachable sets and proving that controlled augmentation is equivalent
  to sampling from a fiber at a specified radius.
\item We introduce multi-axis \emph{semantic budget vectors}
  (\S\ref{sec:budget}) and a \emph{semantic velocity} field that grounds
  curvature constraints in classical differential geometry.
\item We propose \emph{Noether-style semantic invariants} (\S\ref{sec:conservation})
  formalizing conservation of concepts as conserved charges under admissible
  transformations.
\item We prove a \emph{Budget Allocation Principle} for \emph{semantic colored
  noise} (\S\ref{sec:colorednoise}), showing that completion-weighted perturbation
  maximizes diversity subject to a fixed budget, in direct analogy with colored
  noise diffusion sampling \cite{davidson2026colorednoisediffusionsampling}.
\item We model mutation profiles as \emph{profile-dependent metric tensors}
  with a metric-learning interpretation (\S\ref{sec:profiles}).
\end{enumerate}


% ════════════════════════════════════════════════════════════════════════════
\section{Background}
\label{sec:background}
% ════════════════════════════════════════════════════════════════════════════

\subsection{Document Augmentation}

Synonym substitution using lexical resources such as WordNet
\cite{miller1995wordnet} was among the first systematic approaches to document
augmentation.  Back-translation \cite{sennrich2016backtranslation} introduced
a noise-injection mechanism via round-tripping through an intermediate language.
Large pre-trained language models \cite{raffel2020t5, brown2020gpt3} enable
fluent paraphrase generation but without quantitative control over semantic
drift.  Automatic evaluation metrics such as BERTScore \cite{zhang2020bertscore}
and BLEURT \cite{sellam2020bleurt} provide post-hoc assessment but do not
constrain generation.

\subsection{Semantic Embedding Spaces}

Dense document representations \cite{mikolov2013efficient, devlin2019bert,
reimers2019sentencebert} map text into high-dimensional spaces where cosine
similarity approximates semantic relatedness.  A substantial literature
\cite{ethayarajh2019contextual, rahaman2019spectralbias} has characterized the
geometric structure of these spaces, identifying directions corresponding to
sentiment, topic, formality, and other document properties.  Information
geometry \cite{amari2016information} provides a rigorous framework for treating
such spaces as Riemannian manifolds equipped with the Fisher--Rao metric;
our profile-dependent metric tensors are a discrete analogue.

\subsection{Diffusion Processes and Colored Noise}

Denoising diffusion probabilistic models \cite{ho2020ddpm} frame generation as
iterative noise removal along a stochastic trajectory.  Recent work has
characterized the spectral bias of these processes: low-frequency (coarse)
structure resolves early in the denoising trajectory, while high-frequency
(fine) detail resolves late \cite{rahaman2019spectralbias,
falck2025fourierspace}.  \emph{Colored noise diffusion sampling}
\cite{davidson2026colorednoisediffusionsampling} exploits this structure by
routing perturbation energy toward unresolved frequency bands via a completion
field, rather than injecting white noise uniformly.  Blue-noise approaches
\cite{huang2024bluenoise} and geometric analyses of sampling trajectories
\cite{wang2025rotationaltrajectories} further develop the idea that
perturbation energy should respect the existing structure of the space.

\subsection{Constraint-Based Optimization and Control Theory}

Constrained text generation has been approached through lexical constraints
\cite{hokamp2017lexically}, Lagrangian relaxation, and retrieval-augmented
consistency enforcement.  In control theory, constrained trajectory planning
on smooth manifolds \cite{lee2013smooth, docarmo1992riemannian} is a mature
field; our budget constraints and curvature bounds are direct analogues of
energy bounds and kinematic constraints in that literature.  Optimal control
theory \cite{liberzon2012optimal, bertsekas1995dynamic} provides the
variational framework underlying our budget allocation principle.


% ════════════════════════════════════════════════════════════════════════════
\section{The Semantic Manifold and Equivalence Classes}
\label{sec:manifold}
% ════════════════════════════════════════════════════════════════════════════

Let $\calX$ denote the space of all syntactically valid documents in a given
natural language.  Under a fixed embedding function $\phi : \calX \to \bR^d$
we work in the embedded space $\phi(\calX)$, which we treat as a subset of
$\bR^d$ with sufficient regularity to support differentiation and distance
computation.

\begin{definition}[Semantic embedding]
\label{def:embedding}
A \emph{semantic embedding} is a map $\phi : \calX \to \bR^d$ such that
$\norm{\phi(D)} = 1$ for all $D \in \calX$ and $\phi(D_1) \approx \phi(D_2)$
whenever $D_1$ and $D_2$ are judged semantically equivalent by human annotators.
\end{definition}

\begin{definition}[Semantic manifold]
\label{def:manifold}
The \emph{semantic manifold} $\calM \subset \bR^d$ is the image $\phi(\calX)$,
endowed with the metric inherited from $\bR^d$ restricted to $\calM$.
\end{definition}

The semantic manifold hypothesis asserts that $\calM$ has non-trivial geometric
structure: it is not the full sphere $S^{d-1}$, but a curved submanifold
\cite{lee2013smooth, docarmo1992riemannian} whose topology reflects the
constraints of natural-language meaning.

\subsection{Semantic Equivalence Classes}

Before introducing fibers it is necessary to establish the equivalence relation
that fibers will formalize.

\begin{definition}[Approximate semantic equivalence]
\label{def:equiv}
Two documents $D_1, D_2 \in \calX$ are \emph{semantically $\ve$-equivalent},
written $D_1 \sim_\ve D_2$, if
\[
   \dM\bigl(\phi(D_1),\, \phi(D_2)\bigr) \;\le\; \ve.
\]
\end{definition}

\begin{proposition}[$\sim_\ve$ is an approximate equivalence relation]
\label{prop:equiv}
For any $\ve > 0$, the relation $\sim_\ve$ is reflexive and symmetric.  It is
transitive up to $2\ve$: if $D_1 \sim_\ve D_2$ and $D_2 \sim_\ve D_3$
then $D_1 \sim_{2\ve} D_3$.
\end{proposition}

\begin{proof}
Reflexivity and symmetry follow immediately from the definition and the
symmetry of $\dM$.  Transitivity up to $2\ve$ follows from the triangle
inequality: $\dM(\phi(D_1), \phi(D_3)) \le \dM(\phi(D_1), \phi(D_2)) +
\dM(\phi(D_2), \phi(D_3)) \le 2\ve$.
\end{proof}

The equivalence classes under $\sim_\ve$ are the $\ve$-neighbourhoods of each
point on $\calM$.  Documents in the same class are, by Definition
\ref{def:equiv}, acceptable paraphrases of one another at tolerance $\ve$.

\begin{definition}[Semantic neighbourhood]
\label{def:neighbourhood}
The \emph{semantic neighbourhood} of $D$ at radius $r > 0$ is
\[
   \mathcal{N}_r(D) \;=\; \bigl\{\, D' \in \calX \;\mid\;
   \dM(\phi(D),\, \phi(D')) \le r \,\bigr\}.
\]
\end{definition}

This gives the logical chain: embeddings $\to$ manifold $\to$ equivalence
relation $\to$ neighbourhoods $\to$ fibers $\to$ perturbation.


% ════════════════════════════════════════════════════════════════════════════
\section{Document Perturbation as Fiber Exploration}
\label{sec:fiber}
% ════════════════════════════════════════════════════════════════════════════

\begin{definition}[Semantic projection]
\label{def:projection}
The \emph{semantic projection} is the surjective map
$\vpi : \calX \to \calM$ defined by $\vpi(D) = \phi(D)$.
\end{definition}

\begin{definition}[Semantic fiber]
\label{def:fiber}
The \emph{semantic fiber} over $m \in \calM$ is
\[
   \fiber{m} \;=\; \bigl\{\, D \in \calX \;\mid\; \vpi(D) = m \,\bigr\}.
\]
The approximate fiber at tolerance $\ve$ is
$\fiber{m}_\ve = \{D \in \calX : \dM(\vpi(D), m) \le \ve\}$.
\end{definition}

The fiber $\fiber{m}$ consists of all documents projecting to the same semantic
representation: in the ideal case they are complete paraphrases of one another.
Controlled document augmentation is therefore fiber exploration.

\begin{proposition}[Fiber exploration as constrained augmentation]
\label{prop:fiber}
Augmentation within semantic budget $B$ is equivalent to sampling from the
approximate fiber $\fiber{\vpi(D_0)}_B$ of the original document $D_0$.
\end{proposition}

\begin{proof}
By Definition \ref{def:neighbourhood}, sampling within budget $B$ is sampling
from $\mathcal{N}_B(D_0)$.  By Definition \ref{def:fiber}, the approximate
fiber $\fiber{\vpi(D_0)}_B$ consists of all $D'$ with
$\dM(\phi(D'), \phi(D_0)) \le B$.  Since $\vpi = \phi$, these two sets
coincide.
\end{proof}

\subsection{Reachable Sets}

\begin{definition}[Reachable set]
\label{def:reachable}
The \emph{reachable set} of $D$ under budget $B$ is
\[
   \calR_B(D) \;=\; \bigl\{\, D' \in \calX \;\mid\;
   \dM(\phi(D), \phi(D')) \le B \,\bigr\}.
\]
\end{definition}

A perturbation engine is more \emph{expressive} if it can reach a larger
fraction of $\calR_B(D)$.  This language---standard in control theory
\cite{liberzon2012optimal}---makes the budget interpretation precise: the
budget $B$ determines the volume of accessible semantic space, and
maximizing variation is equivalent to maximizing coverage of $\calR_B(D)$.

\begin{remark}
Proposition \ref{prop:fiber} and Definition \ref{def:reachable} together
show that $\calR_B(D_0) = \fiber{\vpi(D_0)}_B$.  The reachable set \emph{is}
the approximate fiber.  Budget and fiber radius are the same quantity measured
from different perspectives.
\end{remark}

\begin{remark}[Why synonym replacement can fail]
A sequence of individually small synonym-swap steps can exit $\calR_B(D_0)$
if the trajectory curves sharply.  A step of size $\delta$ followed by a
second step of size $\delta$ in a perpendicular direction reaches a point at
distance $\delta\sqrt{2}$ from the origin; two steps in opposite directions
return to the origin.  Bounding total displacement alone does not prevent
this; curvature constraints are required.
\end{remark}


% ════════════════════════════════════════════════════════════════════════════
\section{Semantic Distance Budgets and Trajectories}
\label{sec:budget}
% ════════════════════════════════════════════════════════════════════════════

\begin{definition}[Perturbation operator]
\label{def:operator}
A \emph{perturbation operator} $\calO : \calX \to \calX$ maps a document to a
modified version.  Operators are organized by scale:
\begin{enumerate}[label=(\roman*)]
  \item \emph{Lexical operators}: synonym substitution, register shifts.
        Typical cost $\Delta \approx 0.001$--$0.012$.
  \item \emph{Syntactic operators}: active--passive conversion, clause reordering.
        Typical cost $\Delta \approx 0.008$--$0.020$.
  \item \emph{Semantic operators}: LLM paraphrase, summarization--expansion.
        Typical cost $\Delta \approx 0.015$--$0.040$.
\end{enumerate}
The \emph{cost} of applying $\calO$ to $D$ is
$c(\calO, D) = \dM(\phi(D), \phi(\calO(D)))$.
\end{definition}

\begin{definition}[Semantic budget vector]
\label{def:budget}
A \emph{semantic budget vector} is
$\mathcal{B} = (B_\ell, B_s, B_p, \delta) \in \bR_{>0}^4$
specifying lexical, syntactic, and semantic budgets and a maximum step size
$\delta > 0$.  The total budget is $B = B_\ell + B_s + B_p$.
\end{definition}

\begin{definition}[Admissible perturbation sequence]
\label{def:admissible}
A sequence of operators $(\calO_1, \ldots, \calO_n)$ producing trajectory
$D_0 \to D_1 \to \cdots \to D_n$ is \emph{admissible} with respect to
$\mathcal{B}$ if:
\[
   \dM\bigl(\phi(D_0),\, \phi(D_n)\bigr) \;\le\; B
   \quad \text{(total displacement)}
\]
\[
   \dM\bigl(\phi(D_{k-1}),\, \phi(D_k)\bigr) \;\le\; \delta
   \quad \forall\, k \quad \text{(step bound).}
\]
\end{definition}

The perturbation problem is:
\begin{equation}
\label{eq:opt}
   \max_{(\calO_1,\ldots,\calO_n)\;\text{admissible}} V(D_n)
\end{equation}
where $V : \calX \to \bR$ is a diversity functional, typically taken as
$\dM(\phi(D_0), \phi(D_n))$.

\subsection{Semantic Velocity and Curvature}

We ground curvature constraints in classical differential geometry
\cite{docarmo1992riemannian}.

\begin{definition}[Semantic velocity]
\label{def:velocity}
The \emph{semantic velocity} at step $k$ of a trajectory is
\[
   v_k \;=\; \phi(D_k) - \phi(D_{k-1}) \;\in\; \bR^d.
\]
\end{definition}

\begin{definition}[Discrete trajectory curvature]
\label{def:curvature}
The \emph{discrete curvature} at step $k$ is
\[
   \kappa_k \;=\; 1 - \cos\bigl(\angle(v_k,\, v_{k+1})\bigr)
             \;=\; 1 - \frac{\inner{v_k}{v_{k+1}}}{\norm{v_k}\,\norm{v_{k+1}}}.
\]
$\kappa_k = 0$ corresponds to straight-line (zero-curvature) motion; $\kappa_k
= 2$ corresponds to a reversal.
\end{definition}

With semantic velocity defined, curvature constraints bound the angular change
between successive steps, preventing the trajectory from oscillating or
reversing direction.  A curvature-constrained trajectory is one satisfying
$\kappa_k \le \kappa_{\max}$ for all interior steps $k$.

\begin{proposition}[Curvature reduces effective reachable radius]
\label{prop:curvature}
For a trajectory with $n = B/\delta$ maximum-size steps, bounding curvature
$\kappa_k \le \kappa_{\max} < 1$ reduces the maximum possible displacement
from the $\sqrt{n}\,\delta$ attainable by unconstrained steps.
\end{proposition}

\begin{proof}[Proof sketch]
Unconstrained steps in mutually orthogonal directions yield displacement
$\norm{\sum_k v_k} \le \sqrt{n}\,\delta$ by the Cauchy--Schwarz inequality on
$\bR^d$.  With $\kappa_k \le \kappa_{\max}$, consecutive velocity vectors lie
within a cone of half-angle $\theta = \arccos(1 - \kappa_{\max})$.  By
induction, the $n$-step displacement is bounded by
$\delta \cdot \sum_{k=0}^{n-1} \cos^k\!\theta$, which is strictly less than
$\sqrt{n}\,\delta$ for $\kappa_{\max} < 1$.  The budget constraint $B$ remains
binding for total cost; curvature constrains the geometry of the reachable set
within the fiber.
\end{proof}

Consider the trajectories
\begin{center}
\emph{researcher} $\to$ \emph{investigator} $\to$ \emph{analyst}
\quad versus \quad
\emph{researcher} $\to$ \emph{detective} $\to$ \emph{prophet.}
\end{center}
Both may fit within a generous budget.  The first has low curvature: each step
moves in a consistent semantic direction.  The second has high curvature:
the direction of motion reverses sharply between steps.  The curvature
constraint formalizes the intuition that coherent paraphrase should feel smooth.


% ════════════════════════════════════════════════════════════════════════════
\section{Conservation Laws}
\label{sec:conservation}
% ════════════════════════════════════════════════════════════════════════════

Many documents contain concepts that must not drift regardless of remaining
budget.  We formalize this as conservation in the sense of classical mechanics.

\begin{definition}[Protected concept set]
\label{def:protected}
A \emph{protected concept set} is a collection
$\mathcal{C} = \{c_1, \ldots, c_k\}$,
each $c_i$ represented by $\phi(c_i) \in \bR^d$.
\end{definition}

\begin{definition}[Conservation constraint]
\label{def:conservation}
A perturbation $D \to D'$ satisfies the \emph{conservation constraint} with
respect to $\mathcal{C}$ and tolerance $\ve > 0$ if
\[
   \bigl|\inner{\phi(D')}{\phi(c_i)} - \inner{\phi(D)}{\phi(c_i)}\bigr|
   \;\le\; \ve
   \qquad \forall\, c_i \in \mathcal{C}.
\]
The projection of the document embedding onto each protected concept axis
must change by at most $\ve$.
\end{definition}

\subsection{Noether-Style Semantic Invariants}

The conservation constraint acquires a deeper interpretation through analogy
with Noether's theorem.  In classical mechanics, each continuous symmetry of
a physical system gives rise to a conserved charge.  We propose an analogous
principle for semantic systems.

\begin{definition}[Semantic conserved quantity]
\label{def:noether}
A function $Q : \calX \to \bR$ is a \emph{semantic conserved quantity} if
every admissible perturbation $D \to D'$ satisfies $|Q(D') - Q(D)| \le \ve$
for some declared tolerance $\ve$.
\end{definition}

\begin{example}[Conserved quantities in practice]
The following functions $Q$ play the role of conserved charges:
\begin{enumerate}[label=(\roman*)]
  \item \emph{Thesis projection}: $Q(D) = \inner{\phi(D)}{\phi(\text{central claim})}$.
  \item \emph{Terminology preservation}: $Q(D) = \inner{\phi(D)}{\phi(t)}$ for
        each domain term $t \in \{\text{RSVP, TARTAN, Friedmann, \ldots}\}$.
  \item \emph{Citation preservation}: $Q(D) = $ fraction of source citations
        appearing in $D$.
  \item \emph{Numerical claim conservation}: $Q(D) = $ set of quantitative
        assertions in $D$, measured by entailment score.
\end{enumerate}
\end{example}

Each conserved quantity defines a constraint hyper-surface in $\calM$; the
intersection of all such hyper-surfaces is the \emph{admissible region} of the
fiber.  A perturbation engine operating within this region cannot accidentally
negate the document's central thesis, rename its theoretical framework, or
drop numerical claims, regardless of budget.  The conservation constraints are
\emph{holonomic} in the variational sense: they restrict the configuration space
to a sub-manifold rather than imposing velocity penalties.

\begin{remark}[Limitations of projection-based conservation]
The projection formulation protects the semantic neighborhood of a concept but
cannot guarantee preservation of multi-word propositional structure.  A synonym
engine could preserve word embeddings individually while reversing the polarity
of a negation or inverting a causal claim.  Full propositional conservation
requires clause-level entailment checking and is left as future work.
\end{remark}


% ════════════════════════════════════════════════════════════════════════════
\section{Semantic Colored Noise}
\label{sec:colorednoise}
% ════════════════════════════════════════════════════════════════════════════

Standard document perturbation is implicitly \emph{semantic white noise}: it
allocates perturbation energy uniformly across all semantic components,
regardless of how much each has already been varied.  We show that this
is suboptimal and introduce a completion-weighted allocation principle
motivated by colored noise diffusion sampling
\cite{davidson2026colorednoisediffusionsampling}.

\subsection{Semantic Completion Fields}

Let $\mathcal{C}_{\mathrm{sem}} = \{c_1, \ldots, c_m\}$ be a decomposition of
the document's semantic content into components.  Concretely:
\[
   \mathcal{C}_{\mathrm{sem}} \;=\;
   \{\text{thesis},\; \text{evidence},\; \text{terminology},\;
     \text{style},\; \text{narrative},\; \text{examples}\}.
\]

\begin{definition}[Semantic completion field]
\label{def:completion}
The \emph{semantic completion field} is a function
$\gamma : \mathcal{C}_{\mathrm{sem}} \times [0,1] \to [0,1]$
where $\gamma(c, t)$ measures the fraction of variation budget already spent
on component $c$ by perturbation step $t$.  A component with
$\gamma(c, t) \approx 1$ is saturated; one with $\gamma(c, t) \approx 0$ is
largely unexplored.
\end{definition}

\begin{example}[Completion landscape of a mature document]
A document with a well-developed thesis and established terminology, but
relatively varied style and few illustrative examples, might exhibit:
\[
   \gamma(\text{thesis}, t) = 0.95, \quad
   \gamma(\text{terminology}, t) = 0.92, \quad
   \gamma(\text{style}, t) = 0.40, \quad
   \gamma(\text{examples}, t) = 0.25.
\]
A uniform perturbation engine would waste budget on thesis and terminology;
a completion-weighted engine routes variation toward style and examples.
\end{example}

\subsection{Semantic Colored Noise and the Budget Allocation Principle}

Inspired by the completion-aware noise routing of Davidson et al.\ 
\cite{davidson2026colorednoisediffusionsampling}, who define a completion field
$\gamma(f, t)$ on Fourier frequency bands and route stochastic energy toward
unresolved frequencies, we propose an analogous allocation over semantic
components.

\begin{definition}[Semantic colored noise]
\label{def:colorednoise}
A perturbation strategy produces \emph{semantic colored noise} if, at each step
$t$, the energy allocated to component $c$ is proportional to
$\sqrt{1 - \gamma(c, t)}$:
\[
   E(c, t) \;=\; \mathcal{B} \cdot
   \frac{1 - \gamma(c, t)}{\displaystyle\sum_j \bigl(1 - \gamma(j, t)\bigr)}.
\]
Uniform allocation $E(c, t) = \mathcal{B}/m$ is \emph{semantic white noise}.
\end{definition}

\begin{theorem}[Budget Allocation Principle]
\label{thm:allocation}
Let $V(D') = \sum_c \sigma_c \cdot \dM(\phi_c(D_0), \phi_c(D'))$ be a
diversity functional decomposed over semantic components with weights
$\sigma_c > 0$, where $\phi_c$ denotes the projection of the embedding onto
component $c$.  Under a fixed total budget $\mathcal{B}$ and the constraint
$\sum_c E(c) = \mathcal{B}$, semantic colored noise allocation maximizes $V$
subject to staying within the fiber $\calR_\mathcal{B}(D_0)$.
\end{theorem}

\begin{proof}
The diversity functional is separable over components.  For each component $c$,
the marginal contribution to $V$ from allocating $E(c)$ units of budget is
$\sigma_c \cdot g_c(E(c))$ where $g_c$ is a concave increasing function (e.g.,
diminishing returns: additional budget in a saturated component yields less
variation).  By Lagrangian optimization over the constraint
$\sum_c E(c) = \mathcal{B}$, the optimum satisfies
$\sigma_c g_c'(E^*(c)) = \lambda$ for all $c$, where $\lambda$ is the Lagrange
multiplier.  For $g_c(E) = \sqrt{E(1 - \gamma(c))}$ (variation grows as the
square root of unexplored budget), this yields
$E^*(c) \propto 1 - \gamma(c)$, which is exactly the colored noise allocation
of Definition \ref{def:colorednoise}.
\end{proof}

\begin{remark}[Relation to CNS]
Davidson et al.\ \cite{davidson2026colorednoisediffusionsampling} define a
completion field on Fourier frequency bands $f$ and route energy toward
$\sqrt{1 - \gamma(f, t)}$.  Theorem \ref{thm:allocation} establishes that the
same energy-routing principle is optimal in the semantic domain when frequency
bands are replaced by semantic components.  The mathematical structure is
identical; what changes is the space over which completion is measured.
\end{remark}

\begin{remark}[Semantic spectral bias]
The CNS paper \cite{davidson2026colorednoisediffusionsampling} identifies a
\emph{spectral bias} in diffusion models: coarse structure resolves before fine
structure.  The semantic analogue is \emph{conceptual bias}: core thesis and
terminology are established early and become increasingly rigid, while style,
examples, and surface phrasing remain underexplored.  A perturbation engine that
ignores this structure wastes budget on already-saturated components.  Theorem
\ref{thm:allocation} shows that completion-weighted allocation is not merely
heuristically appealing but provably optimal.
\end{remark}

The colored noise framework generalizes naturally to the fiber.  The approximate
fiber $\calR_\mathcal{B}(D_0)$ is not uniformly explored: some semantic
directions are saturated while others remain largely unvisited.  The completion
field $\gamma$ on the fiber guides the perturbation engine toward underexplored
regions, converting the search from blind random walk to directed exploration.


% ════════════════════════════════════════════════════════════════════════════
\section{Profile-Dependent Metric Tensors}
\label{sec:profiles}
% ════════════════════════════════════════════════════════════════════════════

Different discourse communities define semantic equivalence differently.  An
academic paraphrase may freely substitute Latinate for Germanic vocabulary; a
legal paraphrase of the same sentence may not, because the choice between
\emph{shall} and \emph{must} carries distinct deontic force.  These differences
are not stylistic preferences; they reflect genuinely different metric structures
on $\calM$.

\begin{definition}[Profile metric]
\label{def:profile_metric}
A \emph{profile metric} for community $P$ is a positive semi-definite matrix
$M_P \in \bR^{d \times d}$ inducing a distance
\[
   d_P(a, b) \;=\; \sqrt{(a - b)^\top M_P (a - b)}.
\]
The identity $M_P = I$ recovers the standard Euclidean metric.  A
non-identity $M_P$ amplifies distance along directions corresponding to
features that are semantically significant under profile $P$.
\end{definition}

\begin{example}[Academic metric]
Let $\hat{u}_{\mathrm{form}}$ be the unit vector in the direction
$\bar\phi(\text{formal corpora}) - \bar\phi(\text{casual corpora})$.  Then
\[
   M_{\mathrm{acad}} \;=\; I + (w_{\mathrm{form}} - 1)\,
   \hat{u}_{\mathrm{form}} \hat{u}_{\mathrm{form}}^\top, \qquad
   w_{\mathrm{form}} \gg 1
\]
makes movements toward informal register expensive.  A substitution that shifts
formality incurs far higher cost under $d_{\mathrm{acad}}$ than under $d_I$.
\end{example}

\begin{example}[Legal metric]
Under a legal profile, the deontic modality axis (separating obligation,
permission, and prohibition) carries amplified weight.  Substitutions that
alter deontic force are very costly under $d_{\mathrm{legal}}$ even when
embedding cosine similarity is high.
\end{example}

\begin{proposition}[Profile independence of the perturbation algorithm]
\label{prop:profile_independence}
The perturbation engine solving (\ref{eq:opt}) is profile-independent:
it calls $d_P(\cdot, \cdot)$ as an oracle and runs identically for all
profiles.  Changing profile changes the geometry of the search without
altering the algorithm.
\end{proposition}

This proposition means the framework is extensible to new registers without
modifying the perturbation engine.  A domain expert defines a profile by
specifying which semantic axes matter and by how much; the engine automatically
applies that geometry to all subsequent perturbations.

\subsection{Metric Learning Interpretation}

The profile metric $M_P$ has a natural interpretation as a \emph{learned
Mahalanobis metric} \cite{amari2016information}.  Given pairs of documents
$(D_i, D_j)$ labeled as semantically equivalent under profile $P$ by domain
experts, one can learn $M_P$ by minimizing
\[
   \sum_{(i,j)\;\text{equiv}} d_P(\phi(D_i), \phi(D_j))^2
   \;+\; \lambda \sum_{(i,j)\;\text{distinct}} \max\!\bigl(0,\, \alpha - d_P(\phi(D_i), \phi(D_j))\bigr)
\]
via gradient descent on $M_P$ subject to $M_P \succeq 0$.  This connects the
profile metric directly to the metric-learning literature
\cite{larsen2016autoencoding} and provides a path from hand-specified profiles
to data-driven ones.  In the current implementation, $M_P$ is initialized from
contrastive text pairs; full metric learning is reserved for future work.


% ════════════════════════════════════════════════════════════════════════════
\section{Multi-Scale Perturbation Operators}
\label{sec:operators}
% ════════════════════════════════════════════════════════════════════════════

Each operator scale draws from its own budget axis and produces a characteristic
distribution of semantic costs.

\subsection{Lexical Operators}

Synonym substitution uses WordNet \cite{miller1995wordnet} with
frequency-weighted selection.  The weight exponent
\[
   w_i \;=\; \mathrm{freq}(s_i)^{\alpha}, \qquad \alpha = 1 - 0.85\,t,
\]
where $t \in [0,1]$ is the lexical temperature, interpolates between
frequency-dominated selection at $t = 0$ and near-uniform sampling at $t = 1$.
The per-token replacement probability $p(t) = 0.01 + 0.08\,t$ controls
substitution frequency independently of synonym pool width.  At
$t \ge 0.5$, an LLM backend is invoked first; WordNet serves as fallback.

\subsection{Syntactic Operators}

Syntactic operators restructure sentences without altering propositional content:
active--passive conversion, clause reordering, appositive insertion.  Syntactic
operators incur higher semantic cost ($\Delta \approx 0.008$--$0.020$) and
require semantic validation.  The canonical failure case is causality inversion:
\emph{A caused B} to \emph{B was caused by A} is safe, while
\emph{A happened before B} to \emph{B happened before A} is not, despite
identical surface form.  All syntactic proposals are submitted to the fiber
explorer before commitment.

\subsection{Semantic Operators}

LLM paraphrase operators rewrite at the sentence or paragraph level and are the
most expensive ($\Delta \approx 0.015$--$0.040$ per sentence).  They are invoked
only when remaining semantic budget $B_p - \mathrm{spent}_p$ exceeds a
threshold.  The LLM system prompt encodes profile information, directing the
model toward register-appropriate outputs consistent with the active profile
metric.


% ════════════════════════════════════════════════════════════════════════════
\section{Implementation Architecture}
\label{sec:implementation}
% ════════════════════════════════════════════════════════════════════════════

The \textsc{DocPerturb} prototype implements the framework in approximately 3,000
lines of Python and TypeScript across five architectural layers.

The \emph{embedding layer} wraps a sentence-transformer model
\cite{reimers2019sentencebert} and provides cached document embeddings.
Caching is critical: for a document with 200 candidate substitutions, budget
checking without caching requires 200 forward passes through the encoder.
Sentence-level caching recomputes only the embedding of the modified sentence,
reducing the cost to approximately $n_{\mathrm{sentences}}$ forward passes.

The \emph{semantic budget tracker} maintains the trajectory $D_0 \to D_k$,
computing total displacement from origin and per-step distance for every
proposed mutation.  It exposes a unified \texttt{FiberExplorer.admit()}
interface that checks both the budget constraint and the conservation
constraint before any mutation is committed.

The \emph{constraint engine} implements the conservation constraint of
Definition \ref{def:conservation}, checking that each admitted mutation
preserves projections onto protected concept axes within tolerance $\ve$.
The completion field $\gamma(c, t)$ is tracked incrementally, updating after
each committed mutation.

The \emph{perturbation operators} implement lexical, syntactic, and semantic
mutations.  Each operator proposes a candidate and submits it to the fiber
explorer; if admitted, the explorer commits the mutation and updates trajectory
state and the completion field.

The \emph{search loop} performs greedy optimization of (\ref{eq:opt}).
Budget is allocated to components according to the colored noise principle
of Definition \ref{def:colorednoise} at each step.  Future work will explore
beam search and simulated annealing within the budget constraint, as both
are natural extensions of the constrained trajectory framework
\cite{bertsekas1995dynamic}.


% ════════════════════════════════════════════════════════════════════════════
\section{Experimental Evaluation}
\label{sec:experiments}
% ════════════════════════════════════════════════════════════════════════════

We outline experimental protocols; full empirical results are reserved for a
follow-up study with access to human annotators and standardized benchmark
corpora.

\paragraph{Semantic fidelity versus temperature.}
For each temperature $t \in \{0.1, 0.3, 0.5, 0.7, 0.9\}$ and each of the six
profiles, we generate 100 variants of 50 source documents drawn from diverse
genres.  We measure embedding cosine distance between source and variant, human
judgments of semantic equivalence on a 5-point scale, and preservation of
named entities and key claims.

\paragraph{Diversity versus budget.}
We measure the convex hull volume of generated variants in embedding space as a
function of budget $B$.  A well-behaved perturbation engine should exhibit
monotonically increasing hull volume, with diminishing returns as the fiber
boundary is approached.

\paragraph{Colored noise versus white noise.}
We compare uniform budget allocation against the completion-weighted allocation
of Theorem \ref{thm:allocation}, measuring semantic diversity (hull volume),
conservation constraint satisfaction, and human ratings of naturalness.  The
prediction is that colored noise allocation achieves higher diversity within the
same budget by avoiding redundant perturbation of saturated components.

\paragraph{Retrieval robustness testing.}
We insert perturbed variants into a retrieval corpus and measure whether a dense
retrieval system correctly identifies them as near-duplicates.  Budget $B$
parameterizes ground-truth semantic distance; precision at varying budgets
reveals the system's effective sensitivity threshold.

\paragraph{Deduplication benchmark generation.}
We generate semantically near-duplicate pairs at calibrated embedding distances.
Pairs within $B = 0.05$ are labeled ``equivalent,'' between $0.05$ and $0.15$
``related,'' and beyond $0.15$ ``distinct,'' providing a ground-truth dataset
for evaluating semantic deduplication systems.


% ════════════════════════════════════════════════════════════════════════════
\section{Applications}
\label{sec:applications}
% ════════════════════════════════════════════════════════════════════════════

\emph{Data augmentation.}  Perturbations within a small budget $B$ preserve
class labels while increasing lexical and syntactic diversity; the budget
provides a formal guarantee that the label has not been corrupted.

\emph{Retrieval-augmented generation (RAG) robustness.}  Semantic perturbations
of queries and source documents at calibrated distances test whether RAG systems
are sensitive to paraphrase at the level required for their application.

\emph{Plagiarism detection calibration.}  Perturbations at increasing budgets
generate test cases at known semantic distances, enabling calibration of detection
systems and identification of the budget threshold at which detection fails.

\emph{Adversarial robustness testing.}  Selectively violating conservation
constraints generates adversarial examples: documents differing from the source
only in protected concept projections, testing downstream sensitivity.

\emph{Controlled paraphrase datasets.}  Automatically generated pairs at
calibrated budgets provide a scalable, reproducible alternative to human
paraphrase annotation, with explicit ground-truth semantic distance labels.

\emph{Theory version control.}  The framework generalizes beyond documents to
any object admitting a semantic embedding: scientific theories, software
specifications, policy documents.  A theory evolves admissibly if each revision
stays within a semantic budget relative to the founding formulation, and
conserved quantities ensure core claims are not inadvertently abandoned.


% ════════════════════════════════════════════════════════════════════════════
\section{Limitations and Open Problems}
\label{sec:limitations}
% ════════════════════════════════════════════════════════════════════════════

\paragraph{Metric tensor construction.}
Profile metrics are currently initialized from hand-specified contrastive pairs.
Without learned metrics, profiles effectively collapse to the identity metric
with different budget thresholds.  Rigorous construction requires either labeled
equivalence judgments from domain experts or large-scale contrastive pre-training
on register-annotated corpora.

\paragraph{Computational cost.}
Each proposed mutation requires an embedding forward pass.  For large documents
with hundreds of candidate substitutions, this is expensive.  Sentence-level
caching, approximate nearest-neighbor fiber-membership testing, and batched
mutation proposals are viable mitigations.

\paragraph{Propositional conservation.}
Projection-based conservation cannot detect mutations that reverse propositional
polarity or causal direction while preserving individual word embeddings.  Full
propositional conservation requires semantic role labeling, dependency parsing,
and clause-level entailment checking.

\paragraph{Completion field estimation.}
Theorem \ref{thm:allocation} assumes the completion field $\gamma(c, t)$ can
be measured.  In practice, decomposing a document's semantic content into
well-defined components and estimating their completion from embedding alone is
non-trivial.  Structured semantic representations (knowledge graphs,
frame-semantic parses) would make this estimation more reliable.

\paragraph{Embedding faithfulness.}
The framework rests on the assumption that embedding cosine distance approximates
human semantic distance.  This assumption is known to fail for negation,
presupposition, and compositional structures not well-captured by
sentence-level representations.


% ════════════════════════════════════════════════════════════════════════════
\section{Future Directions}
\label{sec:future}
% ════════════════════════════════════════════════════════════════════════════

\emph{Learned profile metrics.}  Training profile-specific metric tensors from
expert equivalence judgments using Siamese networks or information-geometric
methods \cite{amari2016information} would make the metric formalism genuinely
operational.

\emph{Adaptive budget allocation.}  The colored noise allocation of Theorem
\ref{thm:allocation} is currently computed from a fixed initial completion
field.  An adaptive policy that updates $\gamma(c, t)$ after each committed
mutation and reallocates budget dynamically would better track the actual
exploration frontier.

\emph{Graph-based semantic representation.}  Replacing flat embeddings with
knowledge-graph or frame-semantic representations would enable conservation
constraints that operate on propositions rather than projections, addressing
negation and causality failure modes.

\emph{Multimodal perturbation.}  The budget and constraint framework
generalizes to any modality admitting a semantic embedding: images, audio,
structured data.  Multi-modal documents could be perturbed jointly under a
shared semantic budget.

\emph{Connection to RSVP and projection formalisms.}  The fiber exploration
framing connects naturally to projection-based theories such as the RSVP
field-theoretic framework \cite{flyxion2025rsvp}, in which admissible
configurations satisfy a projection condition analogous to fiber membership.
The conservation constraint has the structure of a Dirichlet boundary condition
in the RSVP formalism.  More precisely, the RSVP accessibility interpretation
of entropy---the volume of admissible future trajectories---maps directly onto
the semantic budget: a document's perturbation budget is its local semantic
entropy, and the colored noise principle is the optimal strategy for allocating
that entropy across semantic dimensions.  This correspondence suggests that
DocPerturb, RSVP, and the CNS framework are all instances of a single abstract
principle: constrained finite-budget perturbation in which energy should be
routed toward unresolved dimensions rather than distributed uniformly.


% ════════════════════════════════════════════════════════════════════════════
\section{Conclusion}
\label{sec:conclusion}
% ════════════════════════════════════════════════════════════════════════════

We have argued that document perturbation is fundamentally a geometric problem.
A document is a point on a semantic manifold; its admissible paraphrases occupy
the same semantic fiber; and controlled augmentation is constrained trajectory
planning within that fiber under a budget of semantic distance.  The text file
is merely one coordinate representation of an underlying geometric object; the
geometric object is primary.

The framework's central contributions are the semantic equivalence class
formulation, the reachable-set interpretation of the budget, the velocity-grounded
curvature constraint that prevents smooth individual steps from composing into
discontinuous trajectories, the Noether-style conservation law formulation that
pins conceptual invariants as conserved charges, the Budget Allocation Principle
proving that completion-weighted perturbation (semantic colored noise) maximizes
diversity subject to a fixed budget, and the profile-dependent metric tensor that
models genuinely different notions of semantic equivalence across discourse
communities.

The synonym replacement machinery that motivated this work is one implementation
detail.  The framework accommodates any operator that maps documents to documents
and admits a semantic cost: synonym swaps, dependency-tree rewrites, LLM
paraphrasers, style transfers, translation loops, or multimodal transformations
all consume the same budget from the same fiber explorer.

More broadly, the framework is not specific to documents.  Diffusion sampling,
semantic version control, theory evolution, software refactoring, and RSVP
accessibility dynamics can all be viewed as finite-budget perturbation processes
in which energy should be routed toward unresolved dimensions rather than
distributed uniformly.  Document perturbation is simply the first concrete
application of this principle.

\newpage

% ════════════════════════════════════════════════════════════════════════════
% Bibliography
% ════════════════════════════════════════════════════════════════════════════

\bibliographystyle{plain}

% ── Inline BibTeX (no separate .bib file needed) ─────────────────────────────
\begin{thebibliography}{99}

% ── Conceptual Foundations ───────────────────────────────────────────────────

\bibitem{lee2013smooth}
Lee, J.~M. (2013).
\newblock \emph{Introduction to Smooth Manifolds}, 2nd ed.
\newblock Springer.

\bibitem{docarmo1992riemannian}
do~Carmo, M.~P. (1992).
\newblock \emph{Riemannian Geometry}.
\newblock Birkh\"{a}user.

\bibitem{amari2016information}
Amari, S. (2016).
\newblock \emph{Information Geometry and Its Applications}.
\newblock Springer.

\bibitem{liberzon2012optimal}
Liberzon, D. (2012).
\newblock \emph{Calculus of Variations and Optimal Control Theory:
  A Concise Introduction}.
\newblock Princeton University Press.

\bibitem{bertsekas1995dynamic}
Bertsekas, D.~P. (1995).
\newblock \emph{Dynamic Programming and Optimal Control}.
\newblock Athena Scientific.

% ── Diffusion, Colored Noise, and Spectral Structure ─────────────────────────

\bibitem{davidson2026colorednoisediffusionsampling}
Davidson, H., Issachar, N., \& Benaim, S. (2026).
\newblock Colored noise diffusion sampling.
\newblock \emph{arXiv preprint arXiv:2605.30332}.

\bibitem{ho2020ddpm}
Ho, J., Jain, A., \& Abbeel, P. (2020).
\newblock Denoising diffusion probabilistic models.
\newblock \emph{Advances in Neural Information Processing Systems}, 33.

\bibitem{rahaman2019spectralbias}
Rahaman, N., et al. (2019).
\newblock On the spectral bias of neural networks.
\newblock \emph{Proceedings of ICML 2019}.

\bibitem{falck2025fourierspace}
Falck, F., et al. (2025).
\newblock A Fourier space perspective on diffusion models.
\newblock \emph{arXiv preprint arXiv:2505.11278}.

\bibitem{wang2025rotationaltrajectories}
Wang, Y., \& Vastola, J. (2025).
\newblock Geometric structure of diffusion sampling trajectories.
\newblock \emph{arXiv preprint}.

\bibitem{huang2024bluenoise}
Huang, X., et al. (2024).
\newblock Blue noise for diffusion models.
\newblock \emph{SIGGRAPH 2024}.

\bibitem{lipman2022flowmatching}
Lipman, Y., et al. (2022).
\newblock Flow matching for generative modeling.
\newblock \emph{arXiv preprint arXiv:2210.02747}.

\bibitem{albergo2022stochastic}
Albergo, M., \& Vanden-Eijnden, E. (2022).
\newblock Building normalizing flows with stochastic interpolants.
\newblock \emph{arXiv preprint arXiv:2209.15571}.

% ── Semantic Representation and Embeddings ───────────────────────────────────

\bibitem{mikolov2013efficient}
Mikolov, T., Chen, K., Corrado, G., \& Dean, J. (2013).
\newblock Efficient estimation of word representations in vector space.
\newblock \emph{arXiv preprint arXiv:1301.3781}.

\bibitem{devlin2019bert}
Devlin, J., Chang, M.-W., Lee, K., \& Toutanova, K. (2019).
\newblock BERT: Pre-training of deep bidirectional transformers for language
  understanding.
\newblock \emph{Proceedings of NAACL 2019}.

\bibitem{reimers2019sentencebert}
Reimers, N., \& Gurevych, I. (2019).
\newblock Sentence-BERT: Sentence embeddings using Siamese BERT-networks.
\newblock \emph{Proceedings of EMNLP 2019}, 3982--3992.

\bibitem{ethayarajh2019contextual}
Ethayarajh, K. (2019).
\newblock How contextual are contextualized word representations?
\newblock \emph{Proceedings of EMNLP 2019}, 55--65.

\bibitem{larsen2016autoencoding}
Larsen, A.~B.~L., et al. (2016).
\newblock Autoencoding beyond pixels using a learned similarity metric.
\newblock \emph{Proceedings of ICML 2016}.

% ── Language Generation, Paraphrasing, and Augmentation ─────────────────────

\bibitem{miller1995wordnet}
Miller, G.~A. (1995).
\newblock WordNet: A lexical database for English.
\newblock \emph{Communications of the ACM}, 38(11), 39--41.

\bibitem{sennrich2016backtranslation}
Sennrich, R., Haddow, B., \& Birch, A. (2016).
\newblock Improving neural machine translation models with monolingual data.
\newblock \emph{Proceedings of ACL 2016}, 86--96.

\bibitem{raffel2020t5}
Raffel, C., et al. (2020).
\newblock Exploring the limits of transfer learning with a unified
  text-to-text transformer.
\newblock \emph{Journal of Machine Learning Research}, 21(140), 1--67.

\bibitem{brown2020gpt3}
Brown, T., et al. (2020).
\newblock Language models are few-shot learners.
\newblock \emph{Advances in Neural Information Processing Systems}, 33.

\bibitem{hokamp2017lexically}
Hokamp, C., \& Liu, Q. (2017).
\newblock Lexically constrained decoding for sequence generation using grid
  beam search.
\newblock \emph{Proceedings of ACL 2017}, 1535--1546.

\bibitem{wieting2017revisiting}
Wieting, J., \& Gimpel, K. (2017).
\newblock Revisiting recurrent networks for paraphrastic sentence embeddings.
\newblock \emph{Proceedings of ACL 2017}, 2078--2088.

% ── Evaluation ───────────────────────────────────────────────────────────────

\bibitem{zhang2020bertscore}
Zhang, T., et al. (2020).
\newblock BERTScore: Evaluating text generation with BERT.
\newblock \emph{Proceedings of ICLR 2020}.

\bibitem{sellam2020bleurt}
Sellam, T., Das, D., \& Parikh, A. (2020).
\newblock BLEURT: Learning robust metrics for text generation.
\newblock \emph{Proceedings of ACL 2020}.

% ── Related Theory ───────────────────────────────────────────────────────────

\bibitem{flyxion2025rsvp}
Flyxion. (2025).
\newblock Axioms for a falling universe: Field-theoretic foundations of the
  RSVP framework.
\newblock Independent Research Monograph.

\end{thebibliography}

\end{document}
